% !TEX root = ../thesis.tex

\chapter*{Conclusion}
\addcontentsline{toc}{chapter}{Conclusion}
\chaptermark{Conclusion}
\markboth{Conclusion}{Conclusion}


Our results show that the \textsc{wt} algorithm can be an effective potential solution
for The Registrar to ensure all freshmen are registered for a year-long HUM 110
section. Our algorithm also contributes to the study of stable bipartite
one-sided matching with incomplete preferences and group fairness constraints.
By formulating the task as a weighted graph matching problem, we are able to
give theoretical guarantees about every student being matched to a HUM section
in a full, fair, intra-group stable solution. Though students may be assigned to
a class section that is ranked lower on their preference list than they would
get in the current system, our experiments show that students still tend to receive one
of their top three preferences for the tradeoffs of group fairness and ensuring
every student is matched.

In our experiments, we notably do not evaluate the \textsc{wt} algorithm against
adversarial instances of \textsc{fcfs} and \textsc{greedy}. For both of those algorithms where the
outcome is heavily dictated by the registration order, there necessarily exists
certain orders of students that would result in a worst-case matching. Finding
those theoretical worst-case orderings and testing performance against those
cases would be a relevant extension to this project.

Additionally, our experiments were performed on partially synthesized data based
on previous years of class scheduling, where the system may have inherent biases
that would affect the performance of the algorithm in an actual implementation
setting. It would also usually be standard practice to test the $\phi$ value in the
Mallows preference synthesis to evaluate how the algorithm performs on different
levels of deviation from the central vote. We determined that due to how many other
variables we were already testing (data sample year, preference list length,
randomness over $\phi$ = 0.4, various algorithms), adding that as another
variable was beyond the scope of this thesis. That being said, it would be helpful future work to
take the algorithm and results presented here and test the performance of
\textsc{wt} in a real setting.

More further work can be taken in many directions due to how many different
fields of computer science and mathematics this project incorporates. The
\textsc{wt} algorithm produces many inter-group blocking pairs in each output; one could
investigate approximation algorithms or heuristic methods of reducing the number
of blocking pairs to achieve inter-group stability more often and compare the
effectiveness of those solutions to \textsc{wt}. We also do not explicitly design our
algorithm around balancing class sizes beyond the enforcement of class
capacities. There could be much more done with the weight function to
attempt to encode the class size balancing as an additional constraint to the
problem. Furthermore, most mechanism design solutions across stable matching and
fair division are evaluated under \textit{strategy-proofness}, or whether an entity in
the model can “game” the system to achieve a better outcome for themself at the
cost of other entities. We do not evaluate the \textsc{wt} algorithm for
strategy-proofness; extending this project to find a strategy-proof solution
would also be a fruitful direction for further study. 

In conclusion, we hope our algorithm helps open doors for future CS seniors to
collaborate with The Registrar or other Reed institutions to come up with
projects that allow students to use their Reed education to make a lasting,
tangible impact on the Reed community. We also wish all current and future Reed
students luck in creating their class schedules for years to come. 

